\begin{figure}[H]
\centering
\resizebox{\textwidth}{!}{%
\begin{tikzpicture}[
    font=\small,
    caja/.style={
        draw,
        rectangle,
        rounded corners=4pt,
        align=center,
        minimum height=1cm,
        text width=3.9cm,
        fill=white,
        inner sep=5pt
    },
    cajaCentral/.style={
        draw,
        rectangle,
        rounded corners=4pt,
        align=center,
        minimum height=1.3cm,
        text width=6.2cm,
        fill=white,
        inner sep=6pt,
        thick
    },
    linea/.style={-{Stealth[length=4mm]}, thick}
]


\node[caja] (e1) at (-8.5, 6.5)
    {Pérdida de autonomía\\ en gestión académica};

\node[caja] (e2) at (-2.8, 6.5)
    {Sobrecarga de consultas\\ al personal administrativo};

\node[caja] (e3) at (2.8, 6.5)
    {Errores en inscripción\\ y plazos vencidos};

\node[caja] (e4) at (8.5, 6.5)
    {Frustración y abandono\\ del proceso};

\node[caja] (e12) at (-5.6, 4.2)
    {Dependencia de ayuda\\ externa para tareas básicas};

\node[caja] (e34) at (5.6, 4.2)
    {Retraso en la gestión\\ académica del estudiante};


\node[cajaCentral] (pc) at (0, 1.6)
    {La interfaz de WebSIS genera barreras\\ de usabilidad que impiden la gestión\\ académica autónoma del estudiante};


\node[caja] (c1) at (-9.5, -2.0)
    {Arquitectura de\\ información deficiente};

\node[caja] (c2) at (-4.9, -2.0)
    {Ausencia de visibilidad\\ del estado del sistema};

\node[caja] (c3) at (0, -2.0)
    {Diseño orientado a lógica\\ técnica, no al usuario};

\node[caja] (c4) at (4.9, -2.0)
    {Falta de control y\\ recuperación ante errores};

\node[caja] (c5) at (9.5, -2.0)
    {Sin soporte en periodos\\ críticos de uso};

\node[caja] (c12) at (-7.3, -4.8)
    {Carga cognitiva\\ elevada en el usuario};


\draw[linea] (c1.north) -- (pc.south west);
\draw[linea] (c2.north) -- ($(pc.south) + (-1.5,0)$);
\draw[linea] (c3.north) -- (pc.south);
\draw[linea] (c4.north) -- ($(pc.south) + (1.5,0)$);
\draw[linea] (c5.north) -- (pc.south east);


\draw[linea] (c12) -- (c1);
\draw[linea] (c12) -- (c2);


\draw[linea] (pc.north west) -- (e12.south);
\draw[linea] (pc.north east) -- (e34.south);


\draw[linea] (e12) -- (e1);
\draw[linea] (e12) -- (e2);
\draw[linea] (e34) -- (e3);
\draw[linea] (e34) -- (e4);

\end{tikzpicture}%
}
\caption{Árbol de problemas del sistema WebSIS -- UMSS}
\label{fig:arbol-problemas}
\end{figure}